% Chart pack -- editable source.
% Compile:  cd exports && pdflatex chart_pack.tex
% Figures are the vector PDFs written by the generator scripts; rerun those to
% refresh a chart, then recompile. Text below is yours to edit freely.
\documentclass[11pt]{article}

% ---- layout switch -------------------------------------------------------
% Mobile (default): narrow pages, each page sized to its own chart, so every
% figure fills a phone screen with no pinching. Comment out \mobiletrue for
% wide landscape pages meant for a desktop screen or printing.
\newif\ifmobile
\mobiletrue
% --------------------------------------------------------------------------

\ifmobile
  \usepackage[paperwidth=6in,paperheight=9in,
              top=0.30in,bottom=0.30in,left=0.30in,right=0.30in]{geometry}
\else
  \usepackage[a4paper,landscape,top=1.1cm,bottom=1.0cm,left=1.4cm,right=1.4cm]{geometry}
\fi
\usepackage{graphicx}
\usepackage{xcolor}
\usepackage[T1]{fontenc}
\usepackage{lmodern}
\usepackage{microtype}
\setlength{\parindent}{0pt}
\pagestyle{empty}
% pages never break by themselves: each macro below ships exactly one page
\ifmobile\setlength{\textheight}{200in}\fi

\definecolor{captiongrey}{HTML}{444444}
\definecolor{sectiongrey}{HTML}{888888}

\newsavebox{\gecfig}
\newsavebox{\gechdr}
\newlength{\gecpad}\setlength{\gecpad}{0.30in}
\newlength{\gecgap}\setlength{\gecgap}{11pt}

% \chartpage{title}{caption}{figure path}  -- one chart per page
\newcommand{\chartpage}[3]{%
  \clearpage
  \sbox{\gecfig}{\includegraphics[width=\linewidth]{#3}}%
  \sbox{\gechdr}{\begin{minipage}{\linewidth}\raggedright
      {\large\bfseries #1\par}%
      \ifx\relax#2\relax\else\vspace{4pt}{\small\color{captiongrey}#2\par}\fi
    \end{minipage}}%
  \ifmobile
    % page height = padding + header + gap + figure + padding
    \pdfpageheight=\dimexpr\gecpad+\ht\gechdr+\dp\gechdr+\gecgap
                          +\ht\gecfig+\dp\gecfig+\gecpad\relax
    \noindent\usebox{\gechdr}\par\vspace{\gecgap}\noindent\usebox{\gecfig}%
  \else
    \noindent\usebox{\gechdr}\vfill
    \begin{center}%
      \includegraphics[width=\linewidth,height=0.80\textheight,keepaspectratio]{#3}%
    \end{center}\vfill
  \fi
}

% \sectionpage{title}{blurb}
\newcommand{\sectionpage}[2]{%
  \clearpage
  \ifmobile\pdfpageheight=4in\fi
  \vspace*{\ifmobile 0.9in\else 0.30\textheight\fi}%
  \begin{center}%
    {\ifmobile\LARGE\else\Huge\fi\bfseries #1\par}%
    \vspace{14pt}%
    {\ifmobile\normalsize\else\large\fi\color{captiongrey}%
     \begin{minipage}{\ifmobile 0.92\else 0.78\fi\linewidth}\centering #2\end{minipage}\par}%
  \end{center}%
}

\begin{document}

% ---------------------------------------------------------------- title page
\ifmobile\pdfpageheight=5in\fi
\vspace*{\ifmobile 1.1in\else 0.28\textheight\fi}
\begin{center}
  {\ifmobile\LARGE\else\Huge\fi\bfseries AI-compute supply chain: chart pack\par}
  \vspace{18pt}
  \begin{minipage}{\ifmobile 0.92\else 0.80\fi\linewidth}\centering
    \ifmobile\small\else\large\fi\color{captiongrey}
    Current state of the analysis. Data: UN Comtrade + TDM monthly panel (60 HS6 codes, 2017-01 to 2026-04, checked against Atlas) and Atlas annual data. Sections 3 and 4 use the AI-compute codes only. Details: docs/data.md, docs/supply-chain-narrative.md, results/.
  \end{minipage}
  \vspace{22pt}

  {\color{sectiongrey}2026-08-20}
\end{center}

%% -------------------------------------------------------------- 1. The chain as flows
\sectionpage{1. The chain as flows}{2024 trade flows. China and Hong Kong are counted as one (CHK) throughout. Stages 1-5 use annual Atlas data; stages 6-8 our monthly panel.}

\chartpage{The shape of the chain}{Two input branches feed chip making; finished hardware flows out to the right. Pink inputs are used up per chip; orange equipment builds capacity. Blue: the three codes of our monthly panel. Circle size grows with the log of 2024 trade.}{chain_topology.pdf}

\chartpage{The whole chain on one dollar scale}{Band width = 2024 dollars. The chips stage dominates; small input flows are drawn but too thin to see.}{overviews/supply_chain_overview_dollar_coarse_2024.pdf}

\chartpage{The same flows as a country map}{Countries as nodes, flows as edges colored by stage. Shows combined roles: Mexico takes in chips and parts and ships out servers.}{network/supply_chain_network_2024.pdf}

%% -------------------------------------------------------------- 2. The chain in stages
\sectionpage{2. The chain in stages}{One page per stage: exporters on the left, importers on the right, with the model's groupings in between. Groupings fit the 2024 flows closely (R-squared 0.95-0.99).}

\chartpage{Stage 1 -- raw materials}{Reading guide for all stage pages: left = exporters, right = importers, middle = the model's export and import groups. Much raw-materials trade is not chip-related (the codes also cover welding gas, abrasives, fertilizer inputs), so flows scatter widely.}{hubs_nnf/supply_chain_1_raw_materials_nnf_hubs_2024.pdf}

\chartpage{Stage 2 -- wafers and wafer inputs}{Japan dominates the export side.}{hubs_nnf/supply_chain_2_wafers_nnf_hubs_2024.pdf}

\chartpage{Stage 3 -- lithography and optics inputs}{}{hubs_nnf/supply_chain_3_litho_optics_nnf_hubs_2024.pdf}

\chartpage{Stage 4 -- fab equipment}{Machines, not materials: countries buy these to build capacity, and their imports rise 6-12 months before chip exports.}{hubs_nnf/supply_chain_4_equipment_nnf_hubs_2024.pdf}

\chartpage{Stage 5 -- chips}{The biggest stage, and the whole electronics industry -- phones, cars and PCs, not just AI. Most chips flow to Asian assembly; the China sphere absorbs about 43\%.}{hubs_nnf/supply_chain_5_chips_nnf_hubs_2024.pdf}

\chartpage{Stage 6 -- parts and GPU modules (847330)}{}{hubs_nnf/supply_chain_6_parts_nnf_hubs_2024.pdf}

\chartpage{Stage 7 -- baseboards (847180)}{}{hubs_nnf/supply_chain_7_baseboards_nnf_hubs_2024.pdf}

\chartpage{Stage 8 -- AI servers (847150)}{}{hubs_nnf/supply_chain_8_servers_nnf_hubs_2024.pdf}

%% -------------------------------------------------------------- 3. The factor model over time -- AI-compute codes only
\sectionpage{3. The factor model over time -- AI-compute codes only}{This section uses only the three blue codes (parts/GPU modules, baseboards, AI servers), monthly from 2017. Unless the page says otherwise, China and Hong Kong are counted as one. Shaded bands mark the periods between structural breaks.}

\chartpage{How the model works}{The one idea behind every page that follows.}{mfm_schematic.pdf}

\chartpage{When the network changed shape}{The line measures how much the export groupings moved from one month to the next. Flat for six years, then: Aug 2023 (AI demand), Oct 2024 (the quarter memory export controls were signaled), Sep 2025.}{../results/mfm/tvmfm/chn_hkg_bloc/ai_compute/drift.pdf}

\chartpage{Who is in each export group}{At Oct 2024 the China group loses its separate identity and a Singapore-led group appears.}{../results/mfm/tvmfm/chn_hkg_bloc/ai_compute/loadings_export.pdf}

\chartpage{Who is in each import group}{The import side is steady: a US group, a Mexico group, a Taiwan group, and the rest of the world.}{../results/mfm/tvmfm/chn_hkg_bloc/ai_compute/loadings_import.pdf}

\chartpage{Dollars between groups, month by month}{Three of the sixteen channels carry the AI boom: Taiwan to the US, and both legs of the US-Mexico assembly loop. The rest are flat.}{../results/mfm/tvmfm/chn_hkg_bloc/ai_compute/factors.pdf}

\chartpage{Same, with every country separate}{Without merging China and Hong Kong there is one break in nine years: Aug 2023. Demand changed the country-level structure; policy did not.}{../results/mfm/tvmfm/by_country/ai_compute/drift.pdf}

\chartpage{Parts only (847330)}{91 months without a break -- the stable base layer of the network.}{../results/mfm/tvmfm/by_country/847330/drift.pdf}

\chartpage{With both blocs merged}{Merging China+Hong Kong and US+Mexico hides churn inside each bloc and isolates what happens between them. The breaks line up with policy: the 2019 tariff war, the Oct 2022 export controls, the Apr 2025 tariffs.}{../results/mfm/tvmfm/chnhkg_usamex_blocs/ai_compute/drift.pdf}

\chartpage{The tariff signature}{After Apr 2025 the two blocs move together on a single factor -- seen at no other break.}{../results/mfm/tvmfm/chnhkg_usamex_blocs/ai_compute/loadings_export.pdf}

%% -------------------------------------------------------------- 4. Concentration & fragmentation -- AI-compute codes only
\sectionpage{4. Concentration \& fragmentation -- AI-compute codes only}{Monthly measures built on top of the model, for the three blue codes only. Dotted lines mark Jul 2023 and Apr 2025. Interpretation: results/network\_stats/findings.md.}

\chartpage{Two meanings of fragmentation}{Left: do the model's groups overlap (a technical check, little economic signal). Right: the share of trade crossing between the China and US blocs -- down about 70\% since 2017, in two steps (2019, then 2023 on), with no recovery between.}{../results/network_stats/ai_compute/fragmentation.pdf}

\chartpage{Concentration, taken apart}{Overall concentration splits exactly into three parts. The interesting one asks: do concentrated sellers ship to concentrated buyers? Mildly yes, peaking around 2019, fading through the AI era.}{../results/network_stats/ai_compute/decomposition.pdf}

\chartpage{One buyer or many?}{Left: does each export group sell into one import group or several. Right: a data-quality check, elevated during the 2023-24 transition.}{../results/network_stats/ai_compute/linkage.pdf}

\chartpage{Baseboards: the sharpest split}{The cross-bloc share collapses about 20x from its 2021 peak -- the code most exposed to export controls.}{../results/network_stats/847180/fragmentation.pdf}

\chartpage{Parts: still shared}{The cross-bloc share halves but stays the highest of the three codes: the parts trade remains the layer both blocs still share.}{../results/network_stats/847330/fragmentation.pdf}

\end{document}
